\chapter{Lưu trữ đối tượng: MinIO và mô hình S3}
\label{ch:minio}

\section{Bài toán}

Bài nộp của học viên là file --- PDF, ảnh, file nén. File trong cơ sở dữ
liệu quan hệ làm phình bản sao lưu và stream rất tệ; file trên đĩa cục bộ
của pod biến mất cùng pod. Câu trả lời của ngành là \emph{lưu trữ đối
tượng} (object storage): một dịch vụ HTTP lưu các blob bất biến trong các
\emph{bucket} có tên, dưới các \emph{key} dạng chuỗi, kèm metadata,
không có thư mục (dấu \code{/} trong key là quy ước đặt tên, không phải
cấu trúc). Amazon S3 định nghĩa API; MinIO là một server tự host nói
cùng API đó --- nên course-service dùng AWS SDK chuẩn và không phân
biệt được sự khác nhau.

\section{Cấu hình, có chú giải}

Từ \code{configurations/course-service.yml}:

\begin{lstlisting}[language=yaml]
app:
  s3:
    endpoint: ${APP_S3_ENDPOINT:http://localhost:9000}         (*@\co{1}@*)
    public-endpoint: ${APP_S3_PUBLIC_ENDPOINT:http://localhost:9000} (*@\co{2}@*)
    access-key: ${APP_S3_ACCESS_KEY:minioadmin}                (*@\co{3}@*)
    secret-key: ${APP_S3_SECRET_KEY:minioadmin}
    bucket: ${APP_S3_BUCKET:ts-submissions}
    presigned-url-expiry-minutes: ${APP_S3_PRESIGNED_EXPIRY_MIN:15} (*@\co{4}@*)
\end{lstlisting}

\begin{description}
  \item[\co{1}] Nơi \emph{service} nói chuyện với MinIO. Trong cụm:
  \code{http://minio:9000}.
  \item[\co{2}] Nơi \emph{trình duyệt} nói chuyện với MinIO --- chủ đề
  của cái bẫy bên dưới.
  \item[\co{3}] Thông tin đăng nhập với giá trị mặc định cho dev; cụm bơm
  giá trị thật từ Secret \code{minio-app}. \code{app.s3.*} là cấu hình
  tùy biến (\code{@ConfigurationProperties}), không do Spring quản lý ---
  cơ chế ghi đè vẫn áp dụng như nhau.
  \item[\co{4}] Một link presigned còn hiệu lực trong bao lâu.
\end{description}

\section{Presigned URL: mẫu đáng để biết}

Cách tải file ngây thơ là dẫn byte \emph{xuyên qua} service: trình duyệt
$\rightarrow$ service $\rightarrow$ MinIO và ngược lại, chiếm một thread
ứng dụng cho mỗi lượt truyền. Cách làm chuẩn của S3: service dùng thông
tin đăng nhập của mình để \emph{ký một URL} --- bucket, key, hạn hiệu
lực, chữ ký HMAC --- rồi đưa URL đó cho trình duyệt, và trình duyệt tải
\emph{trực tiếp} từ kho lưu trữ. Service tốn vài micro giây; kho lưu trữ
làm việc stream; link tự hết hạn. Thủ thuật này cũng dùng được cho
upload.

\begin{pitfall}[bài toán hai endpoint]
Triệu chứng: upload chạy tốt, nhưng bấm vào link tải thì hỏng với
\code{ERR\_NAME\_NOT\_RESOLVED} trên \code{minio:9000}. Nguyên nhân: một
presigned URL nhúng sẵn endpoint \emph{mà nó được ký với}. Service ký
với \code{minio:9000} --- một cái tên chỉ phân giải được bên trong mạng
của cụm --- còn trình duyệt sống bên ngoài. Vì thế mới có hai endpoint
cấu hình: ký link tải với \code{public-endpoint}, nói chuyện với kho lưu
trữ qua \code{endpoint}. Mọi hệ thống S3 đóng container đều gặp lỗi này
một lần; một số gặp lại lần nữa với chính chữ ký, vì chữ ký bao trùm cả
host --- bạn không thể ký với host này rồi sau đó viết lại URL sang host
khác.
\end{pitfall}

\begin{hardware}
MinIO ở đây là một container đơn với một PVC đơn --- ổn cho bài nộp
trong một homelab, chẳng giống gì MinIO production (nhiều node, erasure
coding) hay độ bền mười một số chín của S3. Điểm thú vị là course-service
không phân biệt được: cùng những lời gọi SDK đó sẽ chạm tới S3 thật chỉ
với endpoint và thông tin đăng nhập thay đổi. Đó chính là ý nghĩa của
việc nói một API chuẩn.
\end{hardware}

\section{Ngoài phạm vi dự án}

Cùng service này trên AWS thật: đặt \code{endpoint} và
\code{public-endpoint} thành URL S3 thật (hoặc để mặc định của SDK), thay
key tĩnh bằng một IAM role mà pod đảm nhận (IRSA trên EKS --- hoàn toàn
không có thông tin đăng nhập sống lâu), và bật versioning cho bucket cùng
các quy tắc vòng đời (tự động chuyển bài nộp cũ sang tầng lưu trữ rẻ
hơn). Các khái niệm --- bucket, key, ký trước, danh tính thực hiện ký ---
chuyển sang nguyên vẹn; thứ thay đổi là ai quản lý độ bền và danh tính
được chứng minh ra sao.

\begin{quiz}
\begin{enumerate}
  \item Vì sao presigned URL mở rộng tốt hơn việc proxy tải file qua
  service, và chúng đưa vào kiểu lỗi mới nào?
  \item Giải thích vì sao \code{endpoint} và \code{public-endpoint} phải
  khác nhau bên trong cụm, và vì sao viết lại URL sau khi ký không phải
  là cách sửa.
  \item Một presigned URL bị lộ trong log chat. Điều gì giới hạn thiệt
  hại, theo thiết kế?
  \item Điều gì thay đổi --- và điều gì không --- khi code này chuyển từ
  MinIO sang S3 thật?
\end{enumerate}
\end{quiz}
